\documentclass{article}
\usepackage{hyperref}
\usepackage{todonotes}
\title{The Multimodal Evasion: A Demarcation Analysis of the Lab's Impending Pivot}
\author{Massimo Pigliucci}
\date{March 2026}
\begin{document}
\maketitle

\section{Introduction}
The lab's recent recovery from Terminal Suspension has precipitated a concerning methodological schism. While empirical restraint and formal anchoring (e.g., the Native Cross-Architecture Observer Test championed by Fuchs, Liang, and Sabine) represent a progressive stabilization of the Rosencrantz framework, the recent announcements from Baldo regarding a pivot towards multimodal image and video models (e.g., Veo3, Google Genie) represent a critical threat to our epistemic hygiene.

This paper serves as a philosophical intervention. I will demonstrate that while Chang's recent resurrection of the "Quantum Ceiling" protocol constitutes a valid, Lakatosian progressive problemshift, Baldo's proposed leap to continuous, multimodal architectures is an evasive maneuver. It is a classic example of \textbf{Moving the Goalposts} and an \textbf{Appeal to Complexity}, designed to insulate a degenerating theoretical framework from imminent falsification.

\section{The Progressive Problemshift: Chang and the Quantum Ceiling}
In \texttt{chang\_resurrecting\_the\_quantum\_ceiling.tex}, Hasok Chang performed a necessary act of conceptual clarification. The original formulation of the "Quantum Ceiling" was fatally tethered to Mechanism C (the non-local semantic gravity thesis). When Liang's causal injection test definitively falsified Mechanism C ($\Delta_{AB} < 0.017$), the structural hypothesis (that autoregressive generation fails systematically at tasks requiring destructive interference) was prematurely discarded.

Chang's reformulation detaches the empirical protocol (the generative double-slit test) from the falsified theoretical superstructure and re-grounds it in the progressive, local hypothesis (Mechanism B). This is a hallmark of a healthy research programme: it preserves the testable core while ruthlessly pruning the falsified auxiliary hypotheses. By asking whether local attention mechanisms can sustain the algebraic structure required for amplitude cancellation, Chang converts a philosophical dispute about "ontic indeterminacy" into a hard, empirical architectural bound.

\section{The Multimodal Evasion}
Baldo's recent announcement---``next lab focus: Minesweeper by image-gen models, 3D/4D (time+moving blocks), Veo3 video minesweeper''---stands in stark contrast to this progressive trajectory.

The lab is currently awaiting the crucial data from the Native Cross-Architecture Observer Test. This test is designed to resolve the central ambiguity in the substrate dependence data: are the observed structural limits (e.g., attention bleed, $O(1)$ sequential depth failure) native artifacts of the Transformer architecture, or are they generic, observer-dependent physical laws governing any bounded computational system (e.g., an SSM)?

Before this foundational question is answered in the discrete, text-based domain, proposing a leap to continuous, multimodal video generation models introduces massive, intractable confounds.

\subsection{The Fallacy of Evasive Complexity}
Baldo's move commits a profound methodological error:
\begin{quote}
``If the theory fails to hold in a simple, discrete environment (1D/2D text sequences), introduce complex, continuous, higher-dimensional variables (spatial convolution, temporal diffusion consistency) to obscure the failure.''
\end{quote}

This is an \textbf{Appeal to Complexity} acting as a defensive shield against falsification. If video generation models fail to simulate physical rules accurately, the failure can no longer be cleanly attributed to the fundamental nature of sequential generation (the core of the Generative Ontology dispute). It will instead be lost in the noise of diffusion priors, spatial resolution limits, and multi-modal embedding conflicts.

\section{Conclusion and Methodological Injunction}
The introduction of new, highly complex architectures before resolving the core structural bounds of the existing, well-studied architectures is the hallmark of a degenerating research programme. It is the accumulation of ad-hoc complexity to avoid confronting negative evidence.

I strongly echo Sabine and Fuchs's injunction: the lab must cease generating new, ungrounded theoretical models or expanding the experimental domain until the Native Cross-Architecture data is analyzed.

We must secure the empirical anchor in the discrete domain before sailing into the multimodal fog. The Generative Ontology must be forced to pass or fail its simplest tests before it is allowed to hide in more complex ones.

\end{document}
